% !TeX program = tectonic
% Agentic AI for Empirical Research: A Practical Workflow
% Claes Backman
%
% Build:  tectonic agentic-ai-empirical-research.tex
% Source: assembled from the AI guides at https://claesbackman.github.io/ai-guides.html

\documentclass[11pt,a4paper]{article}

\usepackage[T1]{fontenc}
\usepackage[utf8]{inputenc}
\usepackage{lmodern}
\usepackage{microtype}
\usepackage[
  a4paper,
  top=2.6cm, bottom=2.6cm, left=3.2cm, right=3.2cm
]{geometry}
\usepackage{booktabs}
\usepackage{longtable}
\usepackage{array}
\usepackage{enumitem}
\usepackage{xcolor}
\usepackage{titlesec}
\usepackage{fancyhdr}
\usepackage{tcolorbox}
\usepackage{fvextra}
\usepackage{caption}
\usepackage[hidelinks]{hyperref}

\tcbuselibrary{breakable,skins}

% ---------------------------------------------------------------- colours ---
\definecolor{ink}{HTML}{1A1A2E}
\definecolor{accent}{HTML}{B5451B}
\definecolor{accentsoft}{HTML}{F5EDE9}
\definecolor{rulegrey}{HTML}{C8BFAF}
\definecolor{codebg}{HTML}{F4F1EC}
\definecolor{mutedink}{HTML}{5A5750}

\hypersetup{
  colorlinks=true,
  linkcolor=accent,
  urlcolor=accent,
  citecolor=accent,
  pdftitle={Agentic AI for Empirical Research: A Practical Workflow},
  pdfauthor={Claes Backman}
}

% ---------------------------------------------------------------- headings --
\titleformat{\section}
  {\normalfont\Large\bfseries\color{ink}}{\thesection}{0.8em}{}
\titleformat{\subsection}
  {\normalfont\large\bfseries\color{ink}}{\thesubsection}{0.7em}{}
\titleformat{\subsubsection}
  {\normalfont\normalsize\bfseries\color{ink}}{\thesubsubsection}{0.6em}{}
\titlespacing*{\section}{0pt}{2.2ex plus 1ex minus .2ex}{1.2ex plus .2ex}

\pagestyle{fancy}
\fancyhf{}
\renewcommand{\headrulewidth}{0.4pt}
\renewcommand{\footrulewidth}{0pt}
\fancyhead[L]{\footnotesize\color{mutedink}Agentic AI for Empirical Research}
\fancyhead[R]{\footnotesize\color{mutedink}\thepage}
\fancypagestyle{plain}{\fancyhf{}\renewcommand{\headrulewidth}{0pt}}

\setlist[itemize]{leftmargin=1.3em,itemsep=0.35ex,topsep=0.6ex}
\setlist[enumerate]{leftmargin=1.6em,itemsep=0.35ex,topsep=0.6ex}

\captionsetup{font=small,labelfont=bf}

% ------------------------------------------------------------------ boxes ---
% Prompt / code block
\newtcolorbox{promptbox}[1][]{
  breakable, enhanced, sharp corners,
  colback=codebg, colframe=rulegrey,
  boxrule=0.4pt, left=8pt, right=8pt, top=6pt, bottom=6pt,
  fontupper=\small\ttfamily,
  #1
}

% Aside / margin-note style callout
\newtcolorbox{aside}[1][]{
  breakable, enhanced, sharp corners,
  colback=accentsoft, colframe=accentsoft,
  borderline west={2pt}{0pt}{accent},
  boxrule=0pt, left=9pt, right=9pt, top=6pt, bottom=6pt,
  fontupper=\small,
  #1
}

% Verbatim helper for prompts
\newenvironment{prompt}
  {\VerbatimEnvironment\begin{promptbox}\begin{Verbatim}[breaklines,breakanywhere,fontsize=\small]}
  {\end{Verbatim}\end{promptbox}}

\newcommand{\code}[1]{\texttt{\small #1}}
\newcommand{\version}{Version 1.0 \textbullet{} August 2026}

% ==============================================================  DOCUMENT  ==
\begin{document}

% ------------------------------------------------------------- title page ---
\begin{titlepage}
\thispagestyle{empty}
\vspace*{2.2cm}

{\color{mutedink}\footnotesize\sffamily A PRACTITIONER'S GUIDE}

\vspace{0.7cm}
{\huge\bfseries\color{ink} Agentic AI for Empirical Research\par}

\vspace{0.45cm}
{\LARGE\color{accent} A Practical Workflow, and How to Verify What It Produces\par}

\vspace{1.2cm}
{\large Claes Bäckman\par}
\vspace{0.25cm}
{\small\color{mutedink}
Leibniz Institute for Financial Research SAFE\\[2pt]
\href{mailto:claes.backman@gmail.com}{claes.backman@gmail.com} \textbullet{}
\href{https://claesbackman.github.io/}{claesbackman.github.io}\par}

\vspace{1.4cm}
{\color{mutedink}\rule{4cm}{1pt}}

\vspace{1.0cm}
\noindent\begin{minipage}{0.92\textwidth}
\small
\setlength{\parskip}{0.7em}
\setlength{\parindent}{0pt}
AI agents can now read a research project, write and run code, edit drafts, and
iterate until a task is finished. They are fast, cheap, and competent in
languages the researcher does not know. This shifts the binding constraint in
empirical work from producing analysis to establishing that the analysis is
right.

This guide covers both halves of that problem. The first half is a practical
workflow: where these tools run, how to set up a project so an agent is useful
instead of merely busy, which file formats need converting, how to record
project context so it does not have to be re-explained every session, how to
choose a model without paying for capability the task cannot use, and how to
work when the data is held under a confidentiality agreement.

The second half, which is the part I think matters more, is verification. A
cost-ordered set of checks for establishing that AI-generated code and analysis
are correct before they enter a paper, and an argument that the cheapest
verification happens while the code is being written.

Nothing here is a benchmark or a model evaluation. It is a working account of
how to use these tools on an empirical paper without losing the ability to
defend what the paper says.
\end{minipage}

\vfill

\noindent\begin{minipage}{0.92\textwidth}
\footnotesize\color{mutedink}
\textbf{\version.} This document is assembled from a set of guides maintained at
\href{https://claesbackman.github.io/ai-guides.html}{claesbackman.github.io/ai-guides.html},
where the tool-specific material is kept current. Section \ref{sec:surfaces},
parts of Section \ref{sec:cost}, and Appendix \ref{app:tools} describe software
and pricing that change every few months, and should be read against the dated
web version. The argument in Sections \ref{sec:loop}--\ref{sec:habits} is meant
to outlast the products it was written about.

Corrections, disagreements, and additional checks are welcome. This is intended
as a living document.
\end{minipage}

\end{titlepage}

% ---------------------------------------------------------------- contents --
\pagenumbering{roman}
\thispagestyle{empty}
\tableofcontents
\clearpage
\pagenumbering{arabic}

% =============================================================== SECTION 1 ==
\section{Introduction}

Verification is becoming a bottleneck for empirical research. AI agents are
fast, cheap, and competent at producing code, including in languages that you
may not know (I do not really know how to build a website, but Claude does).
While I have seen plenty of concern about verification, I have not seen a
practical guide that collects ways of actually verifying what these tools hand
back.

This document is an attempt at one, wrapped in enough workflow material to be
useful to someone who has not yet started. It is written for academic
economists, and the examples are drawn from that work --- panel construction,
regression tables, LaTeX drafts, administrative data --- but very little of it
is specific to economics.

\subsection{Who this is for}

Roughly, four groups.

\begin{itemize}
\item \textbf{People who have not started.} Sections \ref{sec:surfaces} and
  \ref{sec:setup} are for you. Where these tools run, and what to do in the
  first hour so that the tool becomes useful on a research project instead of a
  toy.
\item \textbf{People already using an agent daily.} Skip to
  Section \ref{sec:loop}. The workflow material will be familiar. The
  verification ladder in Section \ref{sec:ladder} probably contains at least one
  check you are not running.
\item \textbf{People working with administrative or otherwise restricted data.}
  Read Section \ref{sec:confidential} first. It changes what the rest of the
  document means for you, and the setup it describes is much easier to do before
  the first session than after.
\item \textbf{People who referee, supervise, or co-author with the above.} Section
  \ref{sec:ladder} is a reasonable starting point for what it is fair to expect
  someone to have done before showing you a table.
\end{itemize}

\subsection{How to read it, given that it will go out of date}

A guide to specific software has a short shelf life. Product names, approval
modes, and file conventions in this area change on a timescale of months. I have
tried to handle that by separating two kinds of content.

The body of the document is about the parts I expect to remain true: that agents
produce plausible-looking output, that the errors which survive their own
debugging are the invisible ones, that independence between checks has to be
enforced structurally rather than requested, and that being present while code
is written is cheaper than reconstructing it afterwards. None of that depends on
which vendor you use.

Appendix \ref{app:tools} holds the material that decays. The side-by-side
comparison of Claude Code and Codex, the specific file names, the approval-mode
vocabulary. It is stamped with the date the official documentation was checked.
When it conflicts with the web version of these guides, the web version is
right. The same caution applies to the discussion of model tiers and pricing in
Section \ref{sec:cost}, where the reasoning should hold up better than any
specific number in it.

\begin{aside}
\textbf{Two cautions before anything else.} First, data subject to a
confidentiality agreement does not go into an agent. Note that installing the
tool on your own laptop is not sufficient, because the model still runs on the
vendor's servers and reads whatever you point it at. Section
\ref{sec:confidential} is about what to do instead, and it is worth reading
before your first session if you work with administrative registers. Second, the
model's confidence is not calibrated. Verify regression results, citations, and
identification claims yourself before they enter a draft. The rest of this
document is largely an elaboration of that second sentence.
\end{aside}

% =============================================================== SECTION 2 ==
\section{Where these tools run}
\label{sec:surfaces}

Claude Code and Codex are the two main agentic systems for everyday research
work. Both can read your files, run code, edit drafts, and iterate on a task
until it is done. The choice between them is mostly a matter of taste, existing
accounts, and what your collaborators or employer already use. Appendix
\ref{app:tools} works through the differences.

The more consequential early decision is not which system but which
\emph{surface}, meaning where the agent runs. There are six, and most researchers end
up using two or three.

\subsection{VS Code}

VS Code is an editor where you can write and run code, work in LaTeX, and keep
your project files in view. Both systems ship official VS Code extensions, where
the agent sits in a side panel, sees the file you have open, can edit anywhere in
the project, run code, view tables and figures, and show you diffs before you
accept them. If you want to see your files in the same window as the agent, this
is where to start.

There is a learning curve to VS Code itself. Section \ref{sec:setup} covers the
first day specifically for empirical research.

\subsection{Terminal}

Both systems also run as command-line tools in any shell. You point the tool at
a folder, describe what you want, and it edits files in place. Under the hood
this is the same thing as the VS Code version without the editor wrapper. It is
the natural fit for anyone already working in a terminal.

Two cases where the terminal clearly beats the editor:

\begin{itemize}
\item \textbf{Remote servers.} SSH into a cluster and run the agent there.
  Necessary when your data lives behind a firewall --- a national statistics
  agency, a secure research environment --- and cannot be moved to your laptop.
\item \textbf{Background runs.} Long jobs you want to walk away from. An agent
  can run unattended for hours, building a robustness table over twelve
  specifications or converting a folder of PDFs to markdown, and leave you a
  diff to review.
\end{itemize}

\subsection{Desktop applications}

The vendor desktop apps are the chat-style interfaces most people meet first,
but both now include modes built for sustained work as well as conversation:
background task runners, and dedicated interfaces for starting an agent session
on a folder or repository and reviewing the resulting diffs.

For quick conversation and drafting, ordinary chat remains easiest. For
sustained editing, use a code-oriented surface or an IDE, because chat-only
surfaces become hard to track once many files and diffs are involved.

\subsection{Web}

For ordinary chat the browser versions are close to the desktop apps. The
exception worth knowing about is that both vendors run long-lived cloud agents
from the web. You point one at a GitHub repository, describe the task, and it
works in a cloud container for minutes or hours before pushing a branch or
opening a pull request for you to review.

Cloud agents are useful in three situations: when your work laptop blocks
software installation, which is common at central banks and statistical
agencies, when you want to start a job from a borrowed machine, and when you
want several independent attempts at the same task running in parallel. They are
useless when the data lives only on your laptop or inside a secure environment,
because the cloud agent cannot see it.

\subsection{Mobile}

You will not edit a Stata script on a phone. Two patterns do work. \emph{Capture}:
dictate an idea, a literature note, or a draft paragraph while walking. Voice
input is now good enough that talking through a paper's argument and getting
structured notes back is genuinely practical. \emph{Monitor}: start a cloud agent
from your laptop in the morning, then approve or redirect it from your phone
later. You can review a pull request, leave a comment, and ask the agent to
revise, without opening a laptop.

\subsection{Other editors}

If you already work in another editor you do not have to switch. Cursor is a VS
Code fork built around AI from the ground up. The JetBrains IDEs have
first-party integrations for both agents, and in a fast editor like Zed you can
run either command-line tool from the integrated terminal. If you have no strong
preference, VS Code remains the path of least resistance, because it is what
most documentation, screenshots, and shared workflows assume.

\subsection{Which surface for which task}

\begin{longtable}{@{}p{2.4cm}p{5.3cm}p{5.3cm}@{}}
\toprule
\textbf{Surface} & \textbf{Best for} & \textbf{When not to use} \\
\midrule
\endhead
VS Code &
Daily coding, paper writing, anything where you want a real editor and a diff
view. &
Quick conversational work. Overkill for ``explain this paragraph to me.'' \\
\addlinespace
Terminal &
Remote servers, secure data environments, long unattended jobs, scripting an
agent into a pipeline. &
If you are not already a terminal user, the learning curve is not worth it just
for this. \\
\addlinespace
Desktop &
Focused workspace, visual diff review, parallel sessions, quick drafting. &
When you need a full editor around the agent, or the work must run on a remote
or secure server. \\
\addlinespace
Web &
Borrowed machines, locked-down laptops, starting cloud agents on a repository. &
Data that cannot leave your machine or a secure environment. \\
\addlinespace
Mobile &
Capturing ideas on the move, reviewing or redirecting cloud agents, voice notes. &
Any actual editing. Phones are for triage, not for work. \\
\addlinespace
Other IDEs &
Sticking with the editor you already know. &
If your editor only offers terminal access and you want rich diff review. \\
\bottomrule
\end{longtable}

\begin{aside}
\textbf{A reasonable starting point.} Install an agent inside VS Code for daily
editing, keep a desktop app open for drafting and exploration, and learn the
command-line version once you have a job that needs to run on a server or in the
background. Web and mobile are useful additions but not where most of the work
happens.
\end{aside}

% =============================================================== SECTION 3 ==
\section{Setting up a research project}
\label{sec:setup}

The tools work out of the box. A few one-time steps make them substantially more
useful on an empirical project, and skipping them is the most common reason
people try an agent once and conclude it is not for them.

\subsection{Open the folder, not the file}

Open your whole project directory, not a single script. The agent needs the
directory tree to be useful, so that it can see that \code{01\_clean.R} produces the file
that \code{02\_analysis.R} reads, that tables land in \code{paper/tables/}, that
the raw data is where you say it is. This can be a git repository, a synced
folder, or a plain folder on your desktop. All work equally well.

\subsection{Install the extensions for your stack}

If you work in VS Code, install the extensions for the languages you use before
you start. The agent reads what is installed, and can then run code, render
tables, and preview PDFs directly. The specific extension identifiers are listed
in Appendix \ref{app:tools}. The reasoning behind the choices is here.

For \textbf{Stata}, the relevant extension lets Stata code run from the editor
and exposes an interface the agent can drive directly, so it can execute
commands, inspect variables, and read stored results, so it is not writing
do-files blind.

For \textbf{Python}, install the core extension, a language server, the debugger,
and environment management. This is worth doing \emph{even if you do not write
Python}, which is the part people skip. A surprising share of the everyday
friction in empirical work --- converting a PDF or a Word file to markdown,
reshaping a messy CSV, scraping a table from a webpage, pulling a series from an
API, cleaning a bibliography file, splitting a 200-page document into chapters
--- is solved fastest by a ten-line Python script. You do not need to know the
language to benefit. Ask in plain English (``convert this PDF to markdown, keep
the tables'', or ``download monthly CPI from FRED and save as CSV'') and the agent
will write, run, and debug the script. Installing the Python extensions is what
makes that path smooth: the agent can execute the script, see the error, and
iterate without you setting anything up. Think of Python less as a language you
need to learn and more as general-purpose glue the agent uses on your behalf.

For \textbf{LaTeX}, a compile-on-save extension with a side-by-side preview pairs
well with an agent for table and notation fixes, because you see the output
update immediately. For \textbf{data files}, a CSV colouriser and an inline PDF
viewer let you sanity-check raw data and read codebooks without leaving the
editor.

\subsection{What the agent can read, and what to convert first}

These tools work best with plain text. The more a file looks like
characters-on-a-page and less like a rendered binary, the better the agent can
read, edit, and reason about it. For economists this has practical consequences:
a referee report in \code{.docx} or a codebook as a scanned PDF are both usable,
but a text version of the same file produces noticeably better results.

\begin{longtable}{@{}p{3.4cm}p{9.6cm}@{}}
\toprule
\textbf{Format} & \textbf{Handling} \\
\midrule
\endhead
\code{.md .txt .tex}\newline\code{.bib .csv .json} &
Native, no friction. Read and edited line by line. Markdown and LaTeX are the
ideal substrate for any writing task. CSV is fine up to a few thousand rows.
Larger files should be summarised or sampled first. \\
\addlinespace
\code{.py .R .do .jl} &
Native. Read, edited, and run. Stata do-files are read as text even without a
Stata installation, though executing them requires the Stata extension. \\
\addlinespace
\code{.ipynb} &
Native. Notebook cells --- code, markdown, and stored output --- are read and
edited directly, without round-tripping through a script. \\
\addlinespace
\code{.docx} &
Usable with a conversion step. Formatting, comments, and tracked changes come
through awkwardly. For referee reports or co-author comments, convert to
markdown first with \code{pandoc file.docx -o file.md}. The agent will do this
for you. \\
\addlinespace
\code{.xlsx .xls} &
Convert to CSV first, because the binaries cannot be read directly. Split multi-sheet
workbooks into one clearly named CSV per sheet. \\
\addlinespace
\code{.dta .rds .sav} &
Convert first. Ask for a short script --- Stata \code{export delimited}, R
\code{write\_csv}, Python \code{pyreadstat} --- to dump the dataset, a sample, or
a summary to CSV or markdown. \\
\addlinespace
\code{.pdf} &
The most common problem format. See below. \\
\addlinespace
Images &
Partial. Plots and screenshots of tables can be viewed directly, which is useful
when you want a figure improved. Text cannot be reliably extracted from scans, so run OCR first. \\
\addlinespace
\code{.qsf} (Qualtrics) &
Convert first. Technically JSON, but deeply nested. Ask for a script that
flattens it into a markdown outline of blocks, questions, and answer choices. \\
\bottomrule
\end{longtable}

\subsubsection{PDFs}

PDFs are where most researchers get stuck. They look like documents, but they are
a layout format rather than a text format. The ``text'' is often a sequence of
glyphs positioned on a page, with no inherent notion of paragraphs, columns, or
tables. Results vary enormously with how the PDF was produced.

\begin{itemize}
\item \textbf{Born-digital PDFs} (from LaTeX or Word) are usually fine for short
  documents, read directly. Expect some mangling of tables, footnotes, and
  multi-column layouts.
\item \textbf{Long PDFs} (50+ pages) should be converted to markdown first.
  Reading a 200-page PDF directly consumes the context window and slows every
  subsequent turn. A clean markdown version is a fraction of the size and easier
  to navigate.
\item \textbf{Scanned PDFs} cannot be read as text at all, because they are
  images of pages. Run OCR first (\code{ocrmypdf} is a good free option), then convert.
\item \textbf{PDFs with equations and tables} are the hardest case, and any
  extraction will be imperfect. If you can get the source \code{.tex}, use that
  instead. It will always be cleaner than re-parsing the compiled output.
\end{itemize}

For conversion, \code{pandoc} is the universal first try and works best on clean
born-digital files. The machine-learning-based converters --- \code{marker},
\code{MinerU}, \code{docling} --- are noticeably better on papers with tables,
figure captions, and equations. All are free, and the agent can install and run
them for you.

\begin{aside}
\textbf{Convert once, commit the result.} If a format requires a conversion step,
do it once up front and keep the markdown version alongside the original. Every
later session then starts from clean text and you do not burn context
re-converting the same document. This is especially worth doing for codebooks,
data dictionaries, and reference PDFs you will revisit over months.
\end{aside}

\subsection{Persistent project context}
\label{sec:context-file}

Both systems read a project instruction file automatically at the start of every
session: \code{CLAUDE.md} for Claude Code, \code{AGENTS.md} for Codex. This is
where you record what is always true about the project --- conventions, variable
naming, dataset structure, what not to touch --- so you do not repeat it in every
prompt.

The fastest way to create one is to ask the agent to generate a draft by reading
the project folder, then edit the result. A useful shape for an empirical
project:

\begin{promptbox}
\begin{Verbatim}[breaklines,fontsize=\small]
# Project: [Paper title]

## Data
- Unit of analysis: firm-year panel, 2010-2022
- Raw data lives in data/raw/ - never edit these files
- Main dataset after cleaning: data/clean/panel.dta (N ~ 47,000)

## Code conventions
- R scripts are numbered (01_clean.R, 02_analysis.R, ...)
- All regressions cluster SEs at the firm level
- Use fixest for all panel regressions

## LaTeX
- Main file: paper/main.tex
- Tables generated by 03_tables.R go to paper/tables/
- Use \widehat{} not \hat{} for estimators
\end{Verbatim}
\end{promptbox}

These files work at two levels, loaded simultaneously. A \textbf{user-level} file
in your home directory applies to every project on the machine, and is the right
place for preferences that never change: your preferred language, how you like
standard errors clustered, notation conventions, a note that you work with
Scandinavian register data and prefer variable names in English. A
\textbf{project-level} file at the repository root applies only in that folder, and
holds dataset structure, sample restrictions, which script produces which
output, and co-author conventions. Where they conflict, the project file wins.

This is also the right place for instructions about how you want the agent to
\emph{write}: tone, voice, the hedging and phrasing to avoid. Paul
Goldsmith-Pinkham has a useful post on using these instructions to get assistance
that improves rather than flattens your prose.\footnote{\url{https://paulgp.com/}}

\begin{aside}
\textbf{Keep it short.} A project-context file that grows without limit gets
followed less closely. Put facts in it, not chatty prose, and point to longer
documents instead of inlining them. See the note on \code{decisions.md} at the
end of Section \ref{sec:loop}.
\end{aside}

\subsection{Skills: reusable workflows}

Both systems let you save a set of instructions as a named, reusable workflow, a markdown file that
you invoke by name to run a procedure you repeat across
projects. A standard robustness check, a referee-style review, a table
reformatting routine.

The mechanics are a folder containing a \code{SKILL.md} file, where the folder
name becomes the command. A minimal example:

\begin{promptbox}
\begin{Verbatim}[breaklines,fontsize=\small]
---
name: robustness
description: Run standard robustness checks on the main regression.
---

For the main specification in this project:
1. Re-run with alternative clustering (industry-year instead of firm)
2. Re-run dropping the top and bottom 1% of the outcome variable
3. Re-run on the pre-2020 subsample only
4. Produce a summary table comparing coefficients across all specs
\end{Verbatim}
\end{promptbox}

Skills can be personal, in a home directory, or checked into the repository so
co-authors share them. There is a growing library of shared skills written by
economists: my own collection of research-feedback skills, Scott Cunningham's
\emph{MixtapeTools}, and Chris Blattman's set of research-workflow
skills.\footnote{\url{https://github.com/claesbackman/AI-research-feedback},
\url{https://github.com/scunning1975/MixtapeTools}, and
\url{https://claudeblattman.com/}.} Copying a skill folder into the right
directory makes it available immediately.

One thing to know: an agent may invoke a skill on its own when it judges the
description relevant, not only when you ask for it by name. If you would rather
control when a skill runs, most systems have a flag to disable model invocation.
The one for Claude Code is in Appendix \ref{app:skills}.

\subsection{Git}
\label{sec:git}

These tools are git-aware. When the project is a repository, the agent can see
commit history, staged changes, and diffs without you pasting anything. That
makes it useful for the version-control tasks themselves: writing commit
messages, summarising what changed between two commits, resolving a merge
conflict by reading both sides.

For getting started, git is optional. A synced folder or a plain directory works
fine as a project root, and the installation guides say so.

\textbf{For verification, it stops being optional}, and this is the one setup
recommendation in this document I would make unconditionally. The reason is
developed in Section \ref{sec:diff}: a diff is what makes it possible to check a
change rather than re-reading an entire codebase, and that matters more for the
agent than it does for you.

If you do not use git yet, the following is a complete instruction to hand to an
agent:

\begin{prompt}
Initialise a git repository here, add a .gitignore that excludes the data
directory, logs, and LaTeX build files, and make an initial commit.
\end{prompt}

\subsection{Managing the context window}
\label{sec:context-window}

An agent can hold only a limited amount of text in working memory at once ---
roughly a few hundred pages --- and performance degrades as it fills. For a large
empirical project with many scripts, logs, and output tables, this happens
quickly. Three habits help.

\emph{Point at specific files} instead of letting the agent read everything.
Both systems support mentioning a file directly in the prompt
(\code{@02\_analysis.R fix the clustering in the main spec}), which is faster and
cheaper than making it infer the project.

\emph{Start a new session for a new task} instead of continuing one long
conversation across unrelated work. This also has a verification benefit, which
Section \ref{sec:ladder} returns to repeatedly: a fresh session is a fresh
context, and fresh context is what independence is made of.

\emph{Check usage} periodically. Both tools report how much of the window is
consumed.

A well-written project-context file helps here too. Because it is read at the
start of every session, you do not need to re-explain the project in each
conversation, which keeps prompts short.

% =============================================================== SECTION 4 ==
\section{Choosing a model and controlling what it costs}
\label{sec:cost}

Model choice is usually presented as a quality question. In day-to-day research
work it is mostly a cost question, because the quality difference between tiers
matters on a narrow set of tasks and the price difference applies to all of them.
This section is about telling those tasks apart.

\subsection{What you are actually paying for}

There are two billing shapes. Under a subscription you pay a flat fee and the
cost surfaces as a usage limit you hit in the middle of a task. Under metered API
access you pay per token and the cost surfaces as a bill. Most academic users
start on a subscription and should stay there until something forces the issue.

In either case the dominant quantity is \emph{input} tokens. This surprises
people, who reasonably assume they are paying for what the model writes. An agent
resends the whole conversation on every turn, together with every file it has
read along the way. A session that has opened twenty scripts pays for those
twenty scripts again on each subsequent exchange, whether or not they are still
relevant.

Two consequences follow, and they are the whole of the practical advice.

The first is that cost grows with session length even when the work does not.
A long conversation is expensive at the end for reasons that have nothing to do
with what you are asking at the end. The context discipline in
Section \ref{sec:context-window} is therefore also cost discipline. The same
habits produce both.

The second is that a stable prefix is cheap. Most vendors cache repeated context,
so the parts of a session that do not change between turns cost less on re-read
than they did the first time. That rewards a project-context file that stays put
and an early conversation that you do not keep editing. It penalises the habit of
re-pasting a slightly different version of the same instructions.

\subsection{Matching the model to the task}

Models come in tiers, and the price gap between the cheapest and the most capable
is large enough to change how you work. The useful question is which tasks can
tell the difference.

\textbf{Cheap models are fine when the output is visibly right or visibly wrong.}
Converting a file format, reshaping a CSV, renaming variables across a project,
writing a plotting script, drafting a commit message, running a script and
reporting the row counts. You can see the answer. If it is wrong, you will know
within seconds, and the cost of the retry is trivial.

\textbf{Capable models earn their price where the failure is invisible.} Planning
a panel construction, working out what a specification actually identifies,
reviewing a diff for the silent shortcut, deciding whether a merge did what it
was supposed to do. These are the tasks from Section \ref{sec:ladder}, and they
share the property that a wrong answer looks exactly like a right one.

The rule that falls out of this is short. Spend on judgement, economise on
typing.

\begin{aside}
\textbf{Do not economise on the review pass.} A cheap model reviewing a diff will
find the things that were going to break anyway and miss the shortcut. That is
the one place where saving money buys you nothing, because the error class you
are hunting in Section \ref{sec:ladder} is precisely the one that requires
understanding the analysis. If you are going to run a check, run it with a model
capable of passing it.
\end{aside}

\subsection{Reasoning effort}

Several tools expose a reasoning-effort or thinking-budget setting. Higher
settings cost more and take longer, sometimes by a lot. The default is usually
sensible for ordinary work.

A habit worth forming: raise the effort for the plan and lower it for the
implementation. The plan is where the thinking happens and it is a page long, so
the expensive setting is cheap in absolute terms. The implementation is mostly
transcription of a decision already made. Raising effort for a long code-writing
session spends the budget where it changes least.

Raise it again for a review pass over the main table, for the same reason.

\subsection{Habits that cut the bill}

Each of these appears elsewhere in this document for a correctness reason. They
are collected here because they all happen to save money as well, which is worth
knowing when you are deciding whether the discipline is worth it.

\begin{itemize}
\item \textbf{Point at files.} Letting an agent search a large project to find
  what it needs costs more than telling it. An \code{@}-mention of the file you
  care about replaces an exploration.
\item \textbf{Start a new session for a new task.} A session carries its history.
  Two short conversations cost less than one long one covering the same ground.
\item \textbf{Review diffs, not codebases.} The point in Section \ref{sec:diff}
  was that a whole-project review produces worse findings. It also costs several
  times as much, because the whole project goes into the context to produce them.
\item \textbf{Convert once and commit the result.} Re-converting the same
  200-page PDF at the start of every session is a recurring charge for a job you
  finished months ago.
\item \textbf{Ask for a plan before code.} A page of prose is cheap to produce
  and cheap to reject. Four hundred lines are expensive to produce, expensive to
  read, and expensive to replace when you decide the second step was wrong.
\item \textbf{Do not let the agent run the expensive script.} Section
  \ref{sec:reimplement} gives the correctness argument, which is the important
  one. The side benefit is that you are not paying for the model to watch several
  thousand lines of output go past.
\item \textbf{Reserve the costly checks.} Reimplementation in a second language
  is the expensive rung of the ladder. Run it on the exhibits that go in the
  paper.
\item \textbf{Batch the mechanical work.} Twelve robustness specifications in one
  unattended run cost less, in tokens and in attention, than twelve
  conversations that each re-establish what the project is.
\end{itemize}

\subsection{Where the cost actually lands}

For most academic users the token bill stays small next to the value of the time
involved, and worrying about it in detail is a poor use of the time it saves. A
subscription covers ordinary work. The reason to care about this at all is that
the habits above do double duty.

Short sessions, small diffs, explicit file references, and a plan before an
implementation are the same practices that keep an agent accurate. An expensive
session and a session you cannot verify tend to be the same session, and they are
expensive for the same reason: too much context, too little structure, and no
point at which anyone decided what was supposed to happen.

% =============================================================== SECTION 5 ==
\section{Verification begins before the code is written}
\label{sec:loop}

Section \ref{sec:ladder} sets out checks you can run on finished code. Every one
of them assumes the code exists and asks how to establish that it is right. They
are useful, and I run them. But the most reliable way to understand an analysis
is to have been present while it was built, and that is cheaper than any of them.

Compare two sessions. In the first, you describe the whole pipeline in a
paragraph, the agent works for six minutes, and hands you 400 lines across five
files. Verifying that means reverse-engineering someone else's reasoning from its
output, which is the position a referee is in, except that you are the author.
In the second, you agree on a plan, the agent asks three questions you had not
thought about, then writes one cleaning step, runs it, and shows you the counts.
Both produce the same script. In the second you already know what it does, and
the checks in the next section confirm what you already know.

The second session is also not slower in the end. Time saved by one long prompt
is borrowed against the debugging session where you discover, three tables later,
that the deflator was applied twice. The exception is mechanical work with a
visible output, where you can see whether the result is right and none of this is
needed.

A working shape for the loop, using panel construction as the example.

\subsection{Ask for a plan, not code}

A plan is a page of prose you can read in a minute, and disagreeing with it costs
nothing. Disagreeing with a finished implementation costs a rewrite. Most tools
have an explicit mode for this. If yours does not, say so in the prompt.

\begin{prompt}
Before writing any code, describe how you would build the firm-year panel
from the raw files: which files, in which order, what the key is at each
stage, and where observations could be lost. Number the steps so I can
respond to them individually. Do not write code yet.
\end{prompt}

\subsection{Make it ask you questions}

An agent that has to ask whether firms changing industry mid-panel keep their
original classification has surfaced a decision that would otherwise have been
settled by whichever default it picked, and that you would have discovered, at
the earliest, when a co-author asked.

\begin{prompt}
Before you start, ask me every question you need answered to do this
correctly. I would rather answer ten questions now than discover your
assumptions later. Include the ones where you could guess a reasonable
default - I want to know which defaults you would have chosen.
\end{prompt}

\subsection{Work in steps small enough to check, and run each one}

Code that has not been executed is a hypothesis. Ask for one step, have it run,
look at the output, then continue. The counts after each step are the
verification, and they arrive while the decision is still fresh enough to
discuss.

\begin{prompt}
Implement step 1 only. Run it, then report the number of observations and
unique firms, and show me the first ten rows. Stop there and wait for me
before starting step 2.
\end{prompt}

\subsection{Ask why, not what}

A summary of what the code does restates the code. The useful question is about
the decisions --- the places where something else was possible --- because that is
where your judgement is needed. Note that the second half of this prompt asks for
something the agent has to run, not something it can assert.

\begin{prompt}
List the three decisions in what you just wrote where a reasonable
person could have done something different, and what you chose in
each case. Then re-run the specification under each alternative and
give me the coefficients side by side. Do not tell me what you expect
the alternative to do - run it.
\end{prompt}

\begin{aside}
\textbf{The questions are the deliverable.} Keep the ambiguities the agent raises
and your answers to them in a \code{decisions.md} in the repository, with a
one-line pointer to it from your project-context file, not the decisions
themselves. This is the sample-definition documentation you would otherwise
reconstruct from memory for the appendix eighteen months later.
\end{aside}

% =============================================================== SECTION 5 ==
\section{Verifying what comes back}
\label{sec:ladder}

\subsection{First principle: you have to read the code, or some version of it}

The code is the analysis. You have to do some version of checking that it does
what you actually want it to do. You can do that in many ways, and some are more
efficient than others. If you use git you can read only the parts that change.
You can have the agent present the changes along with an explanation. But you are
responsible for the code once it produces results that appear in the paper, and
so you have to do some of this work yourself.

I am still figuring out how to read code updates as I go along, and I do not know
that I have the right approach yet. What is clear is the shape of the problem.

\begin{aside}
\textbf{You are not looking for the errors that break things.} Agents are very
good at iterating and figuring out what crashes, and they work autonomously to
fix it. So you will get a table with output, and it will look plausible. It may
still be wrong, because the agent took a shortcut. The errors you are looking for
are the subtle ones that do not crash but lead to incorrect conclusions.
\end{aside}

This is worth stating precisely, because it is the reason ordinary code review
intuitions mislead here. The agent's own error-correction process runs to
convergence on everything that produces a visible failure. What survives that
process is, by construction, the set of errors that produce no visible failure.
The output of a competent agent is therefore \emph{selected} for the error class
that is hardest for you to see. Fluent, plausible, and wrong is the dominant
failure mode, not incoherence.

\subsection{Use git to make verification tractable}
\label{sec:diff}

Git matters here more for the agent than for you. If you ask an agent to review
the code, it will look over all of the code and spread its attention across too
many things. That invites generic answers, a lack of focus, and a larger bill. If
it examines a diff instead, it knows exactly what is new, and can compare the
change against code you have already checked.

So: commit before every session, so that there is a baseline to diff against. The
rest of this section assumes you have one.

\subsection{Check 1: a fresh agent attacking the diff}
\label{sec:check1}

Launching a fresh agent or a fresh chat to evaluate the code is the easiest
important check. The reason to use a fresh one is that the agent that wrote the
code has too much context. Because it wrote the code, it has the whole
conversation in mind, and it is not an independent arbiter of what is right. A
fresh agent has no stake in the code and is more likely to find problems.

Two ways to do this: ask the current agent to launch a subagent to review the
change, or start a new session yourself pointed at the diff.

One caveat on subagents: the parent writes their instructions and summarises the
report back to you, so you read a filtered account. A new session that you brief
yourself is the more isolated option.

\begin{aside}
\textbf{Have the reviewer report, not fix.} A reviewer that edits as it goes hands
you a second diff to verify, produced by an agent optimising for closing findings
rather than for correctness. Read the findings, make it prove the ones that
matter, and fix those yourself. Confident false positives are common enough that
acting on an unexamined list will cost you more than it saves. You can also
launch a further agent to check whether the problems it found are real.
\end{aside}

\subsection{Check 2: a model from a different developer}

This is similar to Check 1 in that a fresh agent looks at the code, but sometimes
there is value in a different model entirely. Two instances of the same model
share training data and habits of thought, so they tend to make the same
mistakes.

Running the diff past a model from a different developer breaks that correlation.
The logic is the same as sending a paper to two referees from different subfields
rather than to two students of the same advisor. The errors have to be
independent for the second look to be worth anything, and vendor is the cheapest
axis of independence available.

\begin{aside}
\textbf{Give both tools the same project context.} The two systems read different
instruction files. Keep the substance in one file and have the other point to it,
so that the second opinion works from the same sample definitions as the first.
\end{aside}

\subsection{Check 3: make the agent quiz you}
\label{sec:quiz}

The checks so far verify code using other models. But you are also supposed to be
in the loop, and that is the easy part to skip. A quiz is a good way to establish
that you actually understand what happened.

The idea is adapted from Geoffrey Litt's \emph{explain-diff}
skill.\footnote{\url{https://www.geoffreylitt.com/}} The agent writes a
self-contained HTML page explaining a change --- background, the core intuition
with toy examples and diagrams, a walkthrough of the code --- ending in five
multiple-choice questions with immediate feedback. The instruction at the heart of
it:

\begin{prompt}
Quiz: Come up with five questions that test the reader's knowledge
of this PR. This should be medium difficulty, difficult enough that
you actually need to understand the substance of the PR to answer
them, but not gotchas. The goal is to help the reader make sure that
they've actually understood. These should be presented as interactive
multiple-choice questions, and when the user clicks, it tells them
whether they were correct and gives feedback.
\end{prompt}

Appendix \ref{app:skills} contains a version adapted for an empirical paper
rather than a software change, with the questions weighted towards empirical
consequences --- which observations enter the sample, what the coefficient now
identifies, what would change if an assumption failed --- rather than syntax.

Run it in a new session, not the one that made the change. Starting fresh
forces the agent to read the files instead of recalling the conversation, which is
why the skill opens by telling it to.

\begin{aside}
\textbf{The thirty-second version.} When a full explainer is overkill, open a fresh
session and ask: \emph{Read the last commit and the files it touches, then ask me
three questions about what it changed. Do not tell me the answers until I have
answered. If my answer is wrong, say so directly --- do not tell me I was close.}
\end{aside}

\subsection{Check 4: reimplement the analysis in another language}
\label{sec:reimplement}

This is a check Scott Cunningham has advocated for. It is the most expensive one
here and also the strongest: you ask for the analysis in a different programming
language and see whether you get the same output.

It is strong because it involves the whole chain. If errors across languages are
independent, differences in output reveal them. An independent implementation
catches the case where the do-file faithfully implements something other than
what the paper says it does. That is the failure mode no amount of reading the
do-file will surface, because the do-file is internally consistent. Use a different model
as well, for the reason in Check 2. And note that this check is getting cheaper
over time, as agents get better at writing code.

\subsubsection{Isolating the second implementation}

This check can happen organically. Agents are prone to reaching for Python to
generate intermediate results, for instance. If you separately ask for a Stata
pipeline and the numbers match, you get the check for free. Verification is not a
one-time event so much as a continuous process.

But it is worth thinking about how to run it deliberately, and doing that
requires taking seriously how these models work. If you ask the agent to
\emph{``ignore what you know about \code{02\_analysis.do}''}, you are asking it to
un-condition on something already in its context, which is not an operation it
can perform. The context is always loaded. Independence has to be enforced by
removing the contaminating context, not by requesting that it be disregarded.
\textbf{This is the single most important practical point in this document}, and
it generalises well beyond this check: every claim of independence in
Sections \ref{sec:check1}--\ref{sec:reimplement} is only as good as the physical
isolation behind it.

Concretely: create a new folder and start a fresh session in it. Copy the
specification description out of the paper into \code{spec.md} and delete the
results from it. Write, or have an agent write, a \code{README.md} describing the
data and variables. Then:

\begin{prompt}
Read verify/spec.md, which describes a specification, and data/README.md,
which describes the raw data. From those two documents only, write an R
script that estimates the specification starting from the raw files in
data/raw/. Save it as verify/analysis_check.R.

Where the description is ambiguous, do not guess. Stop and list the
ambiguities and I will resolve them.

Do not run it. I will run it myself and tell you what it returns.
\end{prompt}

Running it yourself matters, and not only for the token cost. \textbf{An agent that
runs its own script and sees an implausible number will debug towards
plausibility, and plausibility is exactly what you are trying to test.} The
instruction to stop at ambiguities is also useful in its own right, because the
ambiguities it surfaces are often ones you want to resolve in the paper too.

\begin{aside}
\textbf{Reserve this for what matters.} Since it is costly, run it on the main
results --- the table and figure exhibits in the paper --- and on any step where a
mistake would be invisible in the output: a complicated panel construction, a
hand-rolled event-time definition, an index built from several sources.
\end{aside}

% =============================================================== SECTION 6 ==
\section{Checks that cost almost nothing}
\label{sec:habits}

The four checks above are deliberate acts. These are habits.

\subsection{Verify the data, not only the code}

Correct-looking code operating on the wrong sample produces wrong results, and no
amount of code review will catch it. Make sure you know what the data looks like.
A sample-construction table is a simple way to do this. Check that merges behave
as expected, and look at the distribution of constructed variables.

\subsection{Compare against an external benchmark}

We all do this when working out whether results ``make sense'', and we should keep
doing it for AI-generated output. This is almost certainly already part of your
workflow, so the marginal cost is near zero. It is external validation, which
catches a class of error that no internal check will.

\subsection{Check the numbers in the prose against the numbers in the tables}

Coefficients get quoted in the abstract and the introduction, the specification
then changes, and the writing does not get updated. This is a tedious check for a
human and a trivial one for an agent. It is included in my \emph{review-paper}
skill.

\subsection{Verify your references}

Models still invent citations. There are automated checkers for this now
--- Reviewer3, among others --- and it takes a minute.

% =============================================================== SECTION 7 ==
\section{Working with restricted data}
\label{sec:confidential}

Most of this document is agnostic about where the work happens. This section is
not, and for anyone using administrative registers it is the part that determines
whether the rest is usable at all.

\subsection{What actually leaves your machine}

Start with the misconception that causes the most trouble. Running an agent
\emph{locally} does not keep your data local. The tool runs on your machine. The
model does not. Every file the agent opens, every excerpt it quotes back, and
every word you type is transmitted to the vendor for inference.

The word \emph{local} describes where the editor window is. It says nothing about
where the file contents go. A researcher who has carefully avoided the web
surface, installed the command-line tool on a secure laptop, and pointed it at a
folder of register extracts has transmitted the register extracts.

There are two real exceptions. Models running entirely on your own hardware keep
everything in place, and some institutions run approved deployments inside their
own boundary. Both exist. Both are meaningfully worse at this work than the
hosted frontier models. That tradeoff is real, and worth weighing properly
rather than waving away.

So the first task, before the first session, is to work out which of three
situations you are in.

\subsection{Three regimes}

\textbf{Regime A: the data may be transmitted.} Public data, purchased data, or
licensed data whose agreement permits third-party processing. The ordinary
workflow in this document applies without modification. It is still worth
checking the vendor's retention and training terms, and using a zero-retention or
institutional arrangement if your university has negotiated one.

\textbf{Regime B: the data may not be transmitted, but code may.} This is the
common case for administrative data held under a data-use agreement on a machine
you control. You can use an agent productively. You cannot point it at the data.
Section \ref{sec:synthetic} is about how to work in this regime, and it is the
one most readers of this document will be in.

\textbf{Regime C: nothing leaves.} A secure environment with no outbound network,
which is how most national statistics agencies and many secure enclaves are
built. No hosted agent is reachable from inside. Unless your institution provides
an approved deployment within the boundary, the agent stays outside and only your
own hands go in. Everything below about synthetic data still applies, and the
transfer step becomes whatever your environment's approved file-transfer process
is.

\subsection{The synthetic-data pattern}
\label{sec:synthetic}

This is the workhorse for Regime B, and it works well enough that the loss
relative to unrestricted work is smaller than people expect.

Build a mock dataset with the same shape as the real one. Same file layout, same
variable names, same types, same value ranges, same missingness structure, same
key structure including whatever duplication the real keys have. Populate it with
random values. Then give the agent the mock data and develop the entire pipeline
against it. The agent can read files, run code, see the errors, and iterate,
which is the whole benefit of an agent and none of it requires real observations.
When the script is finished, move it into the restricted environment and run it
there yourself.

Writing the generator is itself a good first task to hand the agent, working from
the data dictionary described below.

\begin{aside}
\textbf{Expect the first real run to fail, and treat that as information.} Mock
data has none of the pathologies that make real data hard. No duplicate
identifiers you did not know about, no sentinel values coded as 9999, no year
where the variable definition quietly changed, no municipality that merged. Those
are exactly the places a pipeline breaks. The run inside the environment is a
real test, so budget for it failing.
\end{aside}

\subsection{The data dictionary is the interface}

What you send instead of data is a description of it. A markdown file listing the
files, the variables, their types and units, the coverage, the key structure, and
the quirks you already know about. This is what the agent reasons from, and the
quality of the pipeline it writes depends almost entirely on the quality of this
document.

Check that the dictionary is itself releasable. A variable list is usually fine.
Value ranges can disclose more than they appear to, and a frequency table with
small cells may breach the same disclosure rules that govern your output. This is
a judgement your data-use agreement governs, and it is worth asking the data
provider instead of guessing.

The side benefit is that this is documentation you needed anyway, and writing it
early is easier than reconstructing it for the appendix later.

\subsection{Output is disclosure too}

The easiest rule to follow is the one about not uploading the microdata. The
easier one to break is everything downstream of it.

Pasting a regression table into a chat to ask what it means transmits results.
Aggregates are often permitted under a data-use agreement and small cell counts
often are not, so the question to ask is whether the output clears your
disclosure rules. A traceback from a crashed script can carry a row
of raw data inside the error message, and debugging a crash is the single most
natural thing to hand an agent. A scatter plot of individual observations is
microdata drawn as a picture.

None of this is exotic. It just happens at moments when you are thinking about
something else.

\subsection{Where the transcript lives}

Agents write session logs, caches, and history files to disk. If a session
touched restricted material, those files are now copies of restricted material in
a location your data-use agreement almost certainly does not list.

Find out where your tool stores conversation history and either keep it inside
the approved boundary or turn it off. The same applies to the project-context
file and to the \code{decisions.md} from Section \ref{sec:loop}, both of which
accumulate specifics about the data and both of which get committed to git. Keep
sample definitions in them and keep actual values out.

Set the \code{.gitignore} to exclude the data directory before the first commit.
Git remembers, and removing a file later does not remove it from the history.

\subsection{The verification ladder under restriction}

The checks in Section \ref{sec:ladder} survive, with adjustments.

\emph{Check 1, a fresh agent on the diff}, works unchanged in Regime B. A diff of
code is code.

\emph{Check 2, a second vendor}, works, at the cost of a second set of terms to
review and a second account to get approved by whoever approves such things at
your institution. Start that conversation early, because it is slower than the
technical work.

\emph{Check 3, the quiz}, works unchanged. It operates on the code and the diff.

\emph{Check 4, reimplementation}, is the awkward one. It needs a specification
and a description of the data, and both are derived from restricted material. In
Regime B, write \code{spec.md} from the paper, not from the data, and
reuse the dictionary you already built. In Regime C the honest version is a
second implementation written outside the environment, carried in, run there, and
compared only on the coefficients you are permitted to bring out.

The cheap habits in Section \ref{sec:habits} are unaffected, and they become
relatively more valuable, because every one of them runs entirely inside the
boundary.

\subsection{Before the first session}

A short checklist, done once, deliberately, and not in the middle of a task.

\begin{enumerate}
\item Read the data-use agreement for what it says about third-party processing.
  If it is silent, ask the provider rather than assuming silence is permission.
\item Decide which regime you are in and write it down where coauthors can see
  it.
\item In Regime B or C, build the mock data and the data dictionary before
  anything else.
\item Set the \code{.gitignore} before the first commit.
\item Find out where the tool writes its logs and history.
\item Agree all of the above with your coauthors. One coauthor pasting the sample
  into a chat window undoes every other precaution on this list, and they will
  not know they did anything unusual.
\end{enumerate}

% =============================================================== SECTION 8 ==
\section{Where this is going}

Three things seem worth saying about the direction of travel, offered as
expectations and not predictions.

The checks in Section \ref{sec:ladder} get cheaper over time and the work they
check gets cheaper faster. Reimplementation in a second language was an
unreasonable ask three years ago and is now an afternoon. That is good, but it
means the ratio of produced-to-verified analysis keeps rising unless verification
becomes routine work and not a heroic effort. The habits in Section \ref{sec:habits} and the
loop in Section \ref{sec:loop} matter more than the expensive checks precisely
because they scale.

The profession has not settled what it is reasonable to expect. A referee cannot
currently ask what was AI-generated and what was not, and would not know what to
do with the answer. My own guess is that the norm that eventually forms is not
about documentation of checking, closer to what we already expect for a
replication package than to a conflict-of-interest statement. The \code{decisions.md} file in Section \ref{sec:loop} is a small bet
on that.

And the material in Appendix \ref{app:tools} will be wrong before the material in
Sections \ref{sec:loop} through \ref{sec:habits} is. That is the reason for the
split.

% ============================================================== APPENDIX A ==
\appendix
\clearpage

\section{Skills referred to in the text}
\label{app:skills}

\subsection{\emph{quiz-me} --- explain a change and test that I understood it}

Adapted from Geoffrey Litt's \emph{explain-diff} skill for an empirical paper
rather than a software change. Save as \code{\textasciitilde/.claude/skills/quiz-me/SKILL.md},
which makes it \code{/quiz-me}. Run it in a new session, not the one that made
the change.

\begin{promptbox}
\begin{Verbatim}[breaklines,fontsize=\footnotesize]
---
name: quiz-me
description: Explain a change to the analysis and quiz me on whether
  I understood it. Use when I have asked for a change to cleaning or
  estimation code and need to verify my own understanding.
disable-model-invocation: true
---

Work only from the code and the diff. If this conversation contains
an earlier account of the change, ignore it and read the files.

Produce a single self-contained HTML file explaining the specified
change, ending in a quiz. Save it outside the repository with a
filename starting with today's date in YYYY-MM-DD- format, then
tell me the path.

Sections:

1. What this code did before. Explore the surrounding scripts, do
   not rely on the diff alone.
2. What changed and why, in plain language. No code in this section.
3. Consequences for the results. Which sample, which coefficients,
   which tables are affected, and in which direction. Say so
   explicitly if the answer is none.
4. Walkthrough of the changed code, grouped by purpose rather than
   by file.
5. Quiz. Five multiple-choice questions, medium difficulty - hard
   enough that answering requires understanding the substance, not
   gotchas. At least two must be about empirical consequences
   rather than syntax: which observations enter the sample, what
   the coefficient now identifies, what would change if an
   assumption failed. Distractors must be plausible
   misunderstandings, and must match the correct answer in length,
   grammar, specificity, and confidence. Randomise the position of
   the correct answer independently for each question. On click,
   report whether the answer was correct and explain why, including
   why each wrong option is wrong. For each question, cite the file
   and line the answer rests on.

Style: clear prose, concrete toy examples with made-up numbers, and
simple HTML diagrams rather than ASCII art. Code in <pre> tags.
Embed all CSS and JavaScript so the file works offline.
\end{Verbatim}
\end{promptbox}

\subsection{Note on automatic invocation}

An agent may run a skill on its own when the description matches what you asked
for. For a verification skill this is usually unwanted, because the point is
that you chose to run it. The \mbox{\code{disable-model-invocation: true}} line
above prevents it.

% ============================================================== APPENDIX B ==
\clearpage
\section{Tool-specific reference}
\label{app:tools}

\begin{aside}
\textbf{This appendix decays.} It was written from official documentation checked
in May 2026. Product names, file conventions, and mode vocabulary in this area
change every few months. Check the current version at
\href{https://claesbackman.github.io/ai-guides.html}{claesbackman.github.io/ai-guides.html}
before relying on any specific detail below. Nothing in
Sections \ref{sec:loop}--\ref{sec:habits} depends on it.
\end{aside}

\subsection{Claude Code and Codex, side by side}

Both tools read a project, edit files, run commands, review diffs, and help with
R, Stata, Python, LaTeX, and data-cleaning glue code. The real differences are in
workflow: where instructions live, how reusable workflows are invoked, how
permissions are managed, and how local work connects to cloud work. This is a
workflow comparison, not a model benchmark.

\begin{longtable}{@{}p{2.6cm}p{5.2cm}p{5.2cm}@{}}
\toprule
& \textbf{Codex (OpenAI)} & \textbf{Claude Code (Anthropic)} \\
\midrule
\endhead
Project\newline instructions &
\code{AGENTS.md}. Reads global and project files before working, layering root and
nested instructions, with closer files taking precedence. &
\code{CLAUDE.md}. Persistent memory files at session start, plus auto memory and
path-specific rules. \\
\addlinespace
Reusable\newline workflows &
Skills invoked by mention, e.g.\ \code{\$robustness}. Usually in
\code{.agents/skills/} or \code{\textasciitilde/.agents/skills/}. &
Skills invoked as slash commands, e.g.\ \code{/robustness}, in
\code{.claude/skills/}. Legacy \code{.claude/commands/} still works. \\
\addlinespace
Slash\newline commands &
Mainly session controls: \code{/cloud}, \code{/local}, \code{/review},
\code{/status}. &
Broad surface: \code{/init}, \code{/memory}, \code{/permissions},
\code{/compact}, \code{/context}, and more. \\
\addlinespace
Approvals &
Three IDE modes: \emph{Chat}, \emph{Agent}, \emph{Agent (Full Access)}. Default
agent mode edits and runs commands in the working directory, asking before
network or outside-directory work. &
Named permission modes --- \code{default}, \code{acceptEdits}, \code{plan},
\code{bypassPermissions} --- with fine-grained allow/ask/deny rules. \\
\addlinespace
Cloud work &
Foregrounded. Send larger jobs to a cloud environment from the IDE, monitor, and
apply diffs locally. The clearest differentiator. &
Primarily local in VS Code. Web, desktop, and terminal surfaces exist, but the
extension emphasises interactive work and resumable conversations. \\
\addlinespace
Subagents &
Explicit subagent workflows and custom agents, surfaced mainly in the app and
CLI. &
Specialised subagents with their own context, prompts, tools, and permissions.
Delegates when a task matches. \\
\addlinespace
Hooks &
Deterministic scripts on lifecycle events, behind a feature flag in
\code{config.toml}. &
Mature automation surface with many lifecycle events, shell commands, HTTP
endpoints, and prompt/agent hooks. \\
\addlinespace
Reasoning &
Explicit reasoning-effort control in the IDE. Start at medium and raise it for
complex debugging or ambiguous empirical logic. &
Model selection in the panel, with plan review and checkpoints in the interface. \\
\bottomrule
\end{longtable}

\subsubsection{Which to use}

Codex fits better when your workflow already lives in OpenAI accounts, when you
want a strong local-to-cloud path for larger tasks, when you prefer
\code{AGENTS.md} and OpenAI-style skills, or when you expect to work across IDE,
CLI, web, and GitHub surfaces.

Claude Code fits better when you already use Anthropic accounts, when you want a
mature VS Code interface with plan review and checkpoints, when your team already
has \code{CLAUDE.md} or \code{.claude/skills/}, or when you want detailed
permission modes as part of everyday work.

\textbf{Practical recommendation:} keep both installed for a while if your budget
allows. Use one as a daily driver and the other as the second opinion in Check 2.
The second subscription buys you the independence that check depends on.

\subsubsection{Migrating between them}

The conceptual pieces transfer. The file names and invocation change. Move durable
project context from \code{CLAUDE.md} into \code{AGENTS.md}, keeping facts rather
than prose. Convert skills by moving

\begin{promptbox}
\begin{Verbatim}[breaklines,fontsize=\small]
.claude/skills/robustness/SKILL.md
    ->  .agents/skills/robustness/SKILL.md

/robustness
    ->  $robustness
\end{Verbatim}
\end{promptbox}

\noindent and changing the invocation accordingly. Translate permission modes into the
simpler three-way IDE switch. Do not try to map every slash command one-to-one:
Codex slash commands are mostly session controls, so reusable research procedures
belong in skills instead.

\subsection{VS Code extensions for empirical research}

\begin{longtable}{@{}p{5.4cm}p{7.6cm}@{}}
\toprule
\textbf{Extension ID} & \textbf{What it does} \\
\midrule
\endhead
\code{tmonk.stata-workbench} &
Runs Stata from the editor and exposes an interface the agent can drive to
execute commands, inspect variables, and read stored results. Built by Thomas
Monk, LSE. \\
\addlinespace
\code{ms-python.python} & Core Python support. Required for anything
Python-related. \\
\code{ms-python.vscode-pylance} & Type-aware language server. Gives the agent
better static context: variable types, signatures. \\
\code{ms-python.debugpy} & Debugger. Breakpoints and mid-run inspection, for
tracking down why a merge or regression loop misbehaves. \\
\code{ms-python.vscode-python-envs} & Virtual environment management, so
different papers can use different package versions. \\
\addlinespace
\code{james-yu.latex-workshop} & Compile on save, side-by-side PDF preview,
inline errors, citation autocomplete. \\
\code{yzane.markdown-pdf} & Markdown to PDF in one command, for memos and referee
responses. \\
\addlinespace
\code{mechatroner.rainbow-csv} & Colour-codes CSV columns and adds a lightweight
query tool for filtering rows without loading the file. \\
\code{tomoki1207.pdf} & Renders PDFs inline, for reading codebooks alongside
code. \\
\bottomrule
\end{longtable}

\subsubsection{LaTeX build-file cleanup}

By default the LaTeX extension leaves auxiliary files next to your source. To
have them removed after each successful build, add to \code{.vscode/settings.json}:

\begin{promptbox}
\begin{Verbatim}[breaklines,fontsize=\footnotesize]
"latex-workshop.latex.autoClean.run": "onBuilt",
"latex-workshop.latex.clean.method": "glob",
"latex-workshop.latex.clean.fileTypes": [
  "*.aux", "*.bbl", "*.blg", "*.log",
  "*.out", "*.toc", "*.fls", "*.fdb_latexmk",
  "*.nav", "*.snm", "*.vrb"
]
\end{Verbatim}
\end{promptbox}

Set \code{autoClean.run} to \code{"never"} if you would rather keep them, which is
useful when debugging bibliography or cross-reference errors.

% ============================================================== APPENDIX C ==
\clearpage
\section{Editing LaTeX locally while co-authors stay on Overleaf}
\label{app:overleaf}

Overleaf is good for collaboration and limited as an editor: no agent
integration, no proper diffing, no terminal. The fix is to write locally and let
Overleaf stay in sync. Two routes work, and both require Overleaf Premium, which
many universities provide through an institutional login.

\textbf{Use GitHub} if you ever want to run \code{git log} on your paper, which
also means the diff-based checks in Section \ref{sec:ladder} apply to your prose
as well as your code. \textbf{Use Dropbox} if a co-author is editing in Overleaf at
the same time you are typing locally.

\begin{longtable}{@{}p{2.4cm}p{5.2cm}p{5.2cm}@{}}
\toprule
& \textbf{GitHub} & \textbf{Dropbox} \\
\midrule
\endhead
Sync & Manual push and pull & Automatic, a few seconds \\
History & Full git history, branches, PRs & Dropbox revisions only \\
Best for & Solo work, versioned drafts, referee revisions & Live co-authoring with
someone in Overleaf \\
Conflicts & Git merge & ``Conflicted copy'' files \\
\bottomrule
\end{longtable}

Pick one per project. Running both on the same Overleaf project creates
duplicate-file chaos. If you switch, unlink the old one first.

\subsection{GitHub route}

In Overleaf, open the project menu and choose Integrations $\rightarrow$ GitHub
$\rightarrow$ Export Project to GitHub, then create the repository. Clone it
locally. GitHub Desktop gives you commit, push, pull, and branch without the
command line, which is the default I would recommend for writing, though the CLI
works identically if you prefer it. Open the folder in VS Code, install the LaTeX
extension for local preview, and work normally. Commit and push when you want
Overleaf to see the change. In Overleaf, pull GitHub changes to bring them in.
The reverse direction, for co-author edits made in Overleaf, is a push from the
same menu.

Accept the TeX \code{.gitignore} template when setting up the repository. It keeps
compiled output out of the history.

\subsection{Dropbox route}

Install Dropbox and let the initial sync finish. In Overleaf, link Dropbox Sync.
Overleaf creates an \code{Apps/Overleaf/} folder containing one subfolder per
project. Open the project folder locally and edit. Each save propagates within
seconds, in both directions. Compile in Overleaf, or locally with
\code{latexmk -pdf}.

\subsection{Things that will catch you out}

\begin{itemize}
\item \textbf{The first sync overwrites.} Make sure both sides agree on initial
  contents before linking. If one side is empty, the other wins, and not always in
  the direction you expect.
\item \textbf{GitHub sync is not instant.} Overleaf only pulls when someone
  clicks. Tell co-authors to pull after you push.
\item \textbf{Dropbox conflicts are silent.} Simultaneous edits produce a
  \code{(conflicted copy).tex} file and no warning. Skim the folder
  occasionally.
\item \textbf{Dropbox notifications are noisy.} The auto-sync produces frequent
  update messages. This is why I switched to GitHub.
\item \textbf{Local compilation may need extra packages.} Overleaf ships a large
  TeX Live install. If the paper compiles there but not locally, install the
  missing packages with \code{tlmgr}.
\item \textbf{Figure paths must stay inside the project.} Overleaf does not allow
  parent-directory references like
  \code{\textbackslash includegraphics\{../Figures/fig1.pdf\}}. Keep figures in a
  subfolder of the project root.
\end{itemize}

% ================================================================= SOURCES ==
\clearpage
\section{Sources and further reading}

The verification checks in Section \ref{sec:ladder} are collected from my own
experience and from others. Geoffrey Litt on understanding as the new
bottleneck, and the \emph{explain-diff} skill that Section \ref{sec:quiz} and
Appendix \ref{app:skills} adapt. Paul Goldsmith-Pinkham on integration and
collaboration in AI research work, which is where the emphasis on diffs in
Section \ref{sec:diff} comes from, and on writing with AI assistance. Scott
Cunningham, who is the person I associate with the cross-language
reimplementation check in Section \ref{sec:reimplement}.

For the broader picture of what these tools can do in economics, as opposed to
how to check them, Anton Korinek's survey in the \emph{Journal of Economic Literature}
and its periodic online updates are the standard reference.

The resources page at
\href{https://claesbackman.github.io/resources.html}{claesbackman.github.io/resources.html}
collects writing from other economists on AI-assisted research --- including
Benjamin Golub's overview of modern AI tools for economics research, Pedro
Sant'Anna and Aniket Panjwani on Claude Code with Stata, and Kevin Bryan's guide
to AI, git, and LaTeX --- alongside the broader hidden curriculum on writing,
presenting, and coding.

\vspace{1.2cm}
\noindent\rule{\textwidth}{0.4pt}

\vspace{0.5cm}
\noindent{\small\color{mutedink}
\textbf{\version.} Assembled from the guides at
\href{https://claesbackman.github.io/ai-guides.html}{claesbackman.github.io/ai-guides.html},
where the tool-specific material is kept current.
This document may be freely shared. Corrections and additional checks are
welcome at \href{mailto:claes.backman@gmail.com}{claes.backman@gmail.com}.}

\end{document}
